% ============================================================
% CAS/Python ενότητες — Κεφάλαιο 6: Διαγωνοποίηση
%
% QR code: qr_7f56b87079ec.png  (→ latex/qrcodes/)
% URL: https://nikmatz.github.io/ELADIC_GR/code/ch06/
% ============================================================

\casactivbox%
  {Maxima: Διαγωνοποίηση — $PDP^{-1}$, Ορθογώνια, Τετραγ. Μορφές}%
  {%
    \label{cas:ch6_diagonalization}
    Δίνονται οι πίνακες
    $A=\bigl[\begin{smallmatrix}5&4\\1&2\end{smallmatrix}\bigr]$,
    $B=\bigl[\begin{smallmatrix}3&1\\0&3\end{smallmatrix}\bigr]$,
    $S=\bigl[\begin{smallmatrix}4&2\\2&1\end{smallmatrix}\bigr]$.
    \begin{enumerate}[label=(\alph*),itemsep=2pt]
      \item Για τον $A$: βρείτε $P=[v_1|v_2]$ και $D=\operatorname{diag}(\lambda_1,\lambda_2)$
        με \texttt{eigenvectors}. Επαληθεύστε $PDP^{-1}=A$.
        Υπολογίστε $A^{10}=PD^{10}P^{-1}$.
      \item Για τον $B$: υπολογίστε \texttt{nullspace}$(B-3I)$.
        Γιατί ο $B$ \emph{δεν} είναι διαγωνοποιήσιμος;
      \item Για τον συμμετρικό $S$: εφαρμόστε Gram-Schmidt στα
        ιδιοδιανύσματα, κατασκευάστε ορθοκανονικό $Q$ και
        επαληθεύστε $Q^\top Q=I$ και $Q\Lambda Q^\top=S$
        (φασματικό θεώρημα).
      \item Υπολογίστε $Q(\mathbf{x})=\mathbf{x}^\top S\mathbf{x}$ και
        ταξινομήστε τη μορφή από τα πρόσημα των ιδιοτιμών.
    \end{enumerate}
    \smallskip
    \textbf{Εντολές Maxima:}
    \texttt{eigenvectors(A)}, \texttt{invert(P)},
    \texttt{transpose(Q)}, \texttt{nullspace(B-3I)},
    \texttt{ratsimp(P . D . P\^{}-1)}
  }%
  {https://nikmatz.github.io/ELADIC_GR/code/ch06/}%
  {qr_7f56b87079ec.png}

\subsection*{Δραστηριότητα Python}

\begin{pybox}{Python: Διαγωνοποίηση — SVD, $e^A$, Τετραγωνικές Μορφές}%
             {https://nikmatz.github.io/ELADIC_GR/code/ch06/}%
             {qr_7f56b87079ec.png}
\label{py:ch6_diagonalization}
\begin{enumerate}[label=(\alph*),itemsep=2pt]
  \item Με \texttt{np.linalg.eig} διαγωνοποιήστε τον $A$ και
    επαληθεύστε $PDP^{-1}=A$. Υπολογίστε $A^{10}$ με $PD^{10}P^{-1}$
    και με \texttt{matrix\_power} — συγκρίνετε.
  \item Για τον $S$ χρησιμοποιήστε \texttt{np.linalg.eigh}
    (εγγυάται ορθοκανονικά ιδιοδιανύσματα). Επαληθεύστε $Q\Lambda Q^\top=S$.
  \item Εκτελέστε SVD: $M=U\Sigma V^\top$ για έναν $3\times 3$ πίνακα.
    Κατασκευάστε προσέγγιση rank-1 και μετρήστε το σφάλμα Frobenius.
  \item Χρησιμοποιήστε \texttt{scipy.linalg.expm} για να υπολογίσετε
    $e^{\pi/2 \cdot A_0}$ όπου $A_0=\bigl[\begin{smallmatrix}0&-1\\1&0\end{smallmatrix}\bigr]$.
    Επαληθεύστε ότι συμπίπτει με $R(\pi/2)$.
  \item \textit{Προαιρετικό:} Σχεδιάστε ισοϋψείς της τετραγωνικής μορφής
    $Q(x,y)=\mathbf{x}^\top S\mathbf{x}$ (Matplotlib \texttt{contour}).
\end{enumerate}
\medskip
\textbf{Βιβλιοθήκες / Εντολές:}
\texttt{np.linalg.eigh()}, \texttt{np.linalg.svd()},
\texttt{scipy.linalg.expm()},
\texttt{np.linalg.matrix\_power()}, \texttt{plt.contour()}
\end{pybox}
